\reading{LTR-13}{Repurchase Agreements and Financing}{STRIKE FIRST}{2}

\begin{lrkeybox}[Learning objectives — official 2026 spine wording]
\begin{enumerate}[label=\textbf{\alph*.}]
\item Describe the mechanics of repurchase agreements (repos) and calculate the settlement for a repo transaction.
\item Discuss common motivations for entering into repos, including their use in cash management and liquidity management.
\item Explain how counterparty risk and liquidity risk can arise through the use of repo transactions.
\item Assess the role of repo transactions in the collapses of Lehman Brothers and Bear Stearns during the 2007-2009 financial crisis.
\item Compare the use of general and special collateral in repo transactions.
\item Identify the characteristics of special spreads and explain the typical behavior of U.S. Treasury special spreads over an auction cycle.
\item Calculate the financing advantage of a bond trading special when used in a repo transaction.
\end{enumerate}
{\footnotesize\color{FRMGrey}Source: Bruce Tuckman and Angel Serrat, \textit{Fixed Income Securities: Tools for Today's Markets}, 3rd Edition, Chapter 12.}
\end{lrkeybox}

\begin{lrtrapbox}[A correction to the reconciliation record]
The initial block map recorded page 135 of the source notes as a \emph{stranded} LTR-13 page, because it carries the heading ``\texttt{LO 72.g: Calculate the financing advantage of a bond trading special}''. A visual read shows the body beneath that heading is \textbf{Covered Interest Parity} --- LTR-16 material. It is not a stranded page to be moved; it is an \textbf{orphan heading} that the following reading's page inherited.

The consequence matters: \textbf{LTR-13 g has no calculation anywhere in either layer}, and it is on the \S5B.10 formula-recall watchlist. Gap-filled below.
\end{lrtrapbox}

% =====================================================================
\lo{LTR-13 a}{Describe the mechanics of repurchase agreements (repos) and calculate the settlement for a repo transaction.}

\begin{lrdefbox}[Repurchase agreement]
A contract in which a security is traded at some initial price with the understanding that the trade will be \textbf{reversed} at some future date at some fixed price.
\end{lrdefbox}

\begin{lrexbox}[Repo settlement]
On 1 May, counterparty A wishes to borrow \$11 million for 31 days. It sells ABC bonds with a \textbf{face value of \$10 million} and a \textbf{market value of \$11 million} to counterparty B, with a contract price of \$11 million to reflect the bond's market value. Concurrently, A agrees to buy back the bond in 31 days at the contract price plus 0.3\% interest (30 basis points). Calculate the settlement amount.
\tcblower
\textbf{Step 1 --- identify the contract price.} \$11 million, the \emph{market} value. The \$10 million face value is not the base.

\textbf{Step 2 --- apply the repo rate on an actual/360 basis.}
\[ \$11{,}000{,}000 \times \left(1 + \frac{0.3\% \times 31}{360}\right) = \$11{,}000{,}000 \times 1.00025833 = \mathbf{\$11{,}002{,}841.67} \]

\textbf{What this means.} The interest earned over 31 days is \$2,841.67 on \$11 million --- the cost of 31 days of secured funding. The bond returns to A; the cash difference is all that remains.
\end{lrexbox}

\begin{lrtrapbox}[Three inputs, three ways to get it wrong]
\begin{itemize}
\item \textbf{Face versus market value.} The contract price is \textbf{\$11m market value}, not \$10m face. Using face gives \$10,002,583.33.
\item \textbf{Day count.} Repo uses \textbf{actual/360}, not 365. Using 365 gives \$11,002,802.74.
\item \textbf{Haircut and accrued interest.} This example has neither. The general GARP form is \[\text{Settlement} = \bigl[(\text{Price} + \text{Accrued interest}) \times (1 - \text{Haircut})\bigr] \times (1 + \text{Repo rate} \times t)\] and \S5B.9 records this as a confirmed recycled template with \textbf{three distractors, one per omitted input}. Check the stem for all three before computing.
\end{itemize}
\end{lrtrapbox}

% =====================================================================
\lo{LTR-13 b}{Discuss common motivations for entering into repos, including their use in cash management and liquidity management.}

\subsection{Borrowers in repos}

\begin{center}
\small
\begin{tabularx}{\textwidth}{@{}p{3.3cm} X@{}}
\toprule
\rowcolor{FRMSoft}\textbf{Motivation} & \textbf{Mechanism} \\
\midrule
\textbf{Cheap funding} & Repos offer \textbf{lower interest rates than unsecured borrowing}, because of the collateral provided. Lenders are willing to provide funds at favourable rates because the loan is collateralized. \\
\addlinespace
\textbf{Bond financing} & Used to obtain cash to finance a \textbf{long} security position until a buyer is found. A firm buying a \$10 million bond intending to sell at a profit can finance the purchase through an overnight repo, using the bond as collateral. If no buyer appears, it can \textbf{roll} the repo. Rolling with the \emph{same} counterparty is a similar process to the initial trade; rolling with a \emph{different} counterparty requires \textbf{unwinding the initial trade and entering a new one}. \\
\addlinespace
\textbf{Liquidity management} & Firms can borrow via equity, long-term debt, and short-term debt, secured or unsecured. Repos provide secured short-term financing but are \textbf{less stable}, because of the short repayment period and susceptibility to market fluctuations. Equity is the \textbf{most stable} option --- no repayment, no required dividends --- but is costly and demands high expected returns. Striking the balance between financing cost and stability \emph{is} liquidity management. \\
\bottomrule
\end{tabularx}
\end{center}

\subsection{Lenders in repos}

From the lender's perspective, repos can be used for \textbf{either investing or financing}, as part of a cash management or financing strategy.

\begin{center}
\small
\begin{tabularx}{\textwidth}{@{}p{4.1cm} X@{}}
\toprule
\rowcolor{FRMSoft}\textbf{Use} & \textbf{Mechanism} \\
\midrule
\textbf{Cash management} \newline (repos as \emph{investment} vehicles) &
Lenders take the reverse repo side to invest surplus cash in low-risk, short-term instruments. \textbf{Money market mutual funds} and municipalities with idle cash are typical users. Investors favour \textbf{overnight} repos for flexibility; \textbf{open repos} with daily renewals are common, allowing investors to reassess the lending decision. \textbf{Longer maturity repos have lower demand.} Investors prioritize high-quality, typically government-issued or guaranteed collateral, with \textbf{haircuts} to account for potential value declines. Transactions involve margining: margin calls require additional collateral in declining markets but allow \textbf{excess collateral withdrawal in rising markets}. \\
\addlinespace
\textbf{Short position financing} \newline (repos as \emph{financing} vehicles) &
Lenders use reverse repos to obtain the desired bond collateral in order to \textbf{short sell} it, betting on price declines. The transaction flow is similar to a repo but with the intention to short and buy back at lower prices. The strategy takes advantage of the expectation of \textbf{rising} interest rates and \textbf{falling} bond prices. \\
\bottomrule
\end{tabularx}
\end{center}

\begin{lrtrapbox}[The lender's two uses point in opposite directions]
\textbf{Cash management} is a \emph{return-seeking} use of spare cash --- the lender wants the rate. \textbf{Short position financing} is a \emph{collateral-seeking} use --- the lender wants the \emph{bond} and will accept a lower rate to get it. That asymmetry is what creates special rates (LTR-13 e). An option describing a specials trade as motivated by yield has the wrong side of the transaction.
\end{lrtrapbox}

% =====================================================================
\lo{LTR-13 c}{Explain how counterparty risk and liquidity risk can arise through the use of repo transactions.}

\begin{center}
\small
\begin{tabularx}{\textwidth}{@{}p{3.4cm} X@{}}
\toprule
\rowcolor{FRMSoft}\textbf{Risk} & \textbf{How it arises, and how it is managed} \\
\midrule
\textbf{Counterparty risk} \newline (credit risk) &
The risk that the \textbf{borrower} --- the party receiving cash --- defaults on repaying the repo loan and interest. \textbf{Mitigated} by the collateral: if the borrower defaults, the lender can sell the collateral to recover the amount owed. Because repos are typically short-term and collateralized, counterparty risk is \textbf{generally less of a concern} than in other lending arrangements. \\
\addlinespace
\textbf{Liquidity risk} \newline (the bigger concern) &
Arises from the potential for the \textbf{value of the collateral} to move adversely during the repo term. Even where the lender is untroubled by the counterparty's creditworthiness, the collateral can lose value or become illiquid. In market turbulence the value can fall sharply and become hard to sell quickly. Managed by: applying \textbf{haircuts}, making \textbf{margin calls}, \textbf{shortening the repo term}, and accepting only \textbf{high-quality collateral}. \\
\bottomrule
\end{tabularx}
\end{center}

\begin{lrtrapbox}[Which risk is the bigger one, and on whom]
The reading is explicit: counterparty risk is \textbf{lower} in repo because of the collateral, and \textbf{liquidity risk is the bigger concern}. A distractor reversing that ranking is the obvious build. Note also the actor: counterparty risk in a repo is the risk of the \textbf{cash borrower} defaulting, from the \emph{lender's} point of view.
\end{lrtrapbox}

% =====================================================================
\lo{LTR-13 d}{Assess the role of repo transactions in the collapses of Lehman Brothers and Bear Stearns during the 2007-2009 financial crisis.}

Before the crisis the repo market was highly liquid, and both borrowers and lenders were comfortable with lower-quality collateral such as corporate bonds and mortgage-backed securities. As the crisis unfolded, lenders grew cautious, \textbf{demanding better collateral and larger haircuts}.

\begin{center}
\small
\begin{tabularx}{\textwidth}{@{}p{2.6cm} X@{}}
\toprule
\rowcolor{FRMSoft}\textbf{Firm} & \textbf{What happened} \\
\midrule
\textbf{Lehman Brothers} &
Lehman relied heavily on repo funding, using \textbf{tri-party repos} with \textbf{JPMorgan Chase} as clearing and tri-party agent. These repos \textbf{matured daily}, leaving Lehman without funding for the rest of the day; to bridge the gap, JPMorgan lent directly to Lehman intraday, \textbf{initially without haircuts on intraday advances}. As the crisis escalated in \textbf{August 2008}, JPMorgan began imposing haircuts on those intraday loans.

\textbf{Lehman's account:} JPMorgan, despite a conflict of interest arising from its dual agent-and-lender role, breached its duty and took advantage of its insider status, draining nearly \textbf{\$14 billion} in collateral in the days before bankruptcy.

\textbf{JPMorgan's account:} it acted in good faith, providing continued funding up to the last day despite Lehman's deteriorating condition.

Ultimately the dispute over collateral and haircuts exacerbated the collapse. \\
\addlinespace
\textbf{Bear Stearns} &
Bear Stearns initially funded with unsecured commercial paper, later switching to repos with high-quality collateral. During the week of \textbf{10 March 2008} it suffered a run resulting from an unwarranted loss of confidence by certain customers, lenders and counterparties, prompted in part by \textbf{market rumours}. Three related consequences:
\begin{enumerate}
\item Prime brokerage clients withdrew cash and unencumbered securities at a rapid and increasing rate.
\item Repo market lenders declined to roll over or renew repo loans, \textbf{even where supported by high-quality collateral such as agency securities}.
\item Counterparties to non-simultaneous settlements of FX trades refused to pay until Bear Stearns paid first.
\end{enumerate}
The result was a rapid flight of capital that could not be survived. \\
\bottomrule
\end{tabularx}
\end{center}

\begin{lrtrapbox}[The two cases fail differently — do not merge them]
\textbf{Lehman} is a story about \textbf{intraday haircuts imposed by its own clearing agent}, with a conflict of interest at the centre. \textbf{Bear Stearns} is a story about a \textbf{three-sided run} where even good collateral could not buy a rollover. The most instructive detail in the Bear Stearns case is that lenders refused to renew \emph{despite} high-quality agency collateral --- an option asserting that better collateral would have saved it misses the point of the case.
\end{lrtrapbox}

% =====================================================================
\lo{LTR-13 e}{Compare the use of general and special collateral in repo transactions.}

\begin{center}
\small
\begin{tabularx}{\textwidth}{@{}p{2.6cm} X X@{}}
\toprule
\rowcolor{FRMSoft}\textbf{Feature} & \textbf{General collateral (GC)} & \textbf{Special collateral} \\
\midrule
\textbf{Focus} & High-quality, \emph{any} acceptable security from broad categories & A \emph{specific} security requested by the lender \\
\textbf{Repo rate} & \textbf{Highest} --- the GC rate & \textbf{Lower} --- the special rate \\
\textbf{Use case} & General investment & Targeted financing strategies \\
\bottomrule
\end{tabularx}
\end{center}

\begin{lrkeybox}[Why the special rate is the lower one]
The rationale behind GC is that \textbf{specific collateral demands create scarcity}, which \emph{lowers} the repo rate. GC trades are favoured by investors precisely because they usually offer the \textbf{highest} repo rates. Lenders in specials trades seek specific securities --- to short a bond, or to fund inventory --- and accept a trade-off: they get the security they want, \textbf{but at a rate below GC}. Special rates vary by security and term, driven by supply and demand, and importantly the supply and demand for the \emph{underlying security} may not align with its suitability as \emph{collateral} in a specials trade.
\end{lrkeybox}

\begin{center}
\begin{tikzpicture}
\begin{axis}[
  width=11.6cm, height=6.2cm,
  xlabel={Term}, ylabel={Repo rate $R$},
  xlabel style={font=\small\sffamily}, ylabel style={font=\small\sffamily},
  xtick=\empty, ytick=\empty,
  xmin=0, xmax=10, ymin=0, ymax=10,
  axis lines=left,
]
\addplot[FRMNavy, very thick, domain=0.4:9, samples=60]{3.6+2.6*ln(1+x)};
\addplot[FRMGold, very thick, domain=0.4:9, samples=60]{2.4+2.6*ln(1+x)};
\addplot[FRMGreen, very thick, domain=0.4:9, samples=60]{1.2+2.6*ln(1+x)};
\node[font=\scriptsize\sffamily\bfseries, fill=white, inner sep=1.5pt, anchor=west, text=FRMNavy] at (axis cs:7.0,9.2) {FFR};
\node[font=\scriptsize\sffamily\bfseries, fill=white, inner sep=1.5pt, anchor=west, text=FRMGold] at (axis cs:7.0,8.0) {GC rate};
\node[font=\scriptsize\sffamily\bfseries, fill=white, inner sep=1.5pt, anchor=west, text=FRMGreen] at (axis cs:7.0,6.8) {Special rate};
\draw[{Stealth[length=4pt]}-{Stealth[length=4pt]}, FRMRed] (axis cs:4.2,7.4) -- (axis cs:4.2,5.0)
  node[midway, right=2pt, font=\scriptsize\sffamily, fill=white, inner sep=1.5pt, text=FRMRed]{special spread};
\end{axis}
\end{tikzpicture}
\figcap{Three curves, one strict ordering: \textbf{FFR $>$ GC $>$ special}. In the United States the GC rate is typically close to, but \emph{below}, the federal funds rate --- and the fed funds-to-GC spread widened to reflect the \textbf{decreasing supply of Treasury collateral}. The special spread is the vertical distance from GC down to the special rate, which is why its floor is zero and its cap is the GC rate.}
\end{center}

\begin{lrtrapbox}[The ordering is the whole objective]
Any statement placing the special rate \emph{above} GC, or GC \emph{above} the fed funds rate, is wrong by construction. The counterintuitive bit is that \textbf{high demand for a specific bond lowers its repo rate}: the lender of cash is compensated in \emph{collateral}, not in yield.
\end{lrtrapbox}

% =====================================================================
\lo{LTR-13 f}{Identify the characteristics of special spreads and explain the typical behavior of U.S. Treasury special spreads over an auction cycle.}

\begin{lrdefbox}[Special spread]
The difference between the \textbf{GC rate} and the \textbf{special rate} for a specific security and term. Special spreads are closely tied to US government Treasury bond auctions and serve as a gauge of market sentiment.
\end{lrdefbox}

\subsection{The auction cycle}

\begin{itemize}
\item Federal government bonds in the US are sold through \textbf{scheduled auctions}. The most recent issue is the \term{on-the-run (OTR)} or current issue; older issues are \term{off-the-run (OFR)}.
\item OTR issues are \textbf{highly liquid}, making them desirable for both long and short positions, and therefore desirable as collateral.
\item OTR special spreads can exhibit \textbf{daily volatility} depending on the collateral used, and fluctuate over time.
\item \textbf{Special spreads for OTR bonds tend to be narrower immediately after an auction, and wider before an auction.} Narrowing occurs after auctions because of the \textbf{surplus supply} of the new OTR security.
\end{itemize}

\begin{lrtrapbox}[The two layers give different reasons for the pre-auction widening — take the Crash layer's]
The Crash layer says spreads widen before auctions \emph{as demand for the soon-to-be-issued OTR bond rises}. The Lecture layer says they widen \emph{because market participants shift to the new OTR security, and hence demand for old ones decreases}. The Lecture layer's reason does not support its own conclusion: falling demand for a security makes it \emph{less} special, which \textbf{narrows} the spread. The mechanism that widens a spread is rising demand for the security being borrowed. Retained with the Crash layer's reasoning; see the corrections record.
\end{lrtrapbox}

\subsection{Limits of special spreads}

\begin{lrkeybox}[Floor, cap, and the post-2009 refinement]
\begin{enumerate}
\item \textbf{Floor: 0\% on the special \emph{rate}.} When a trader short sells the OTR Treasury but \textbf{fails to deliver} on settlement, the trader does not receive cash from the sale and also misses a day's interest on that cash. To satisfy the delivery obligation, the trader could borrow the bond in the overnight repo market and pay a special rate of \textbf{0\%} --- essentially providing free financing in exchange for receiving the bond. At any rate below 0\% no trader would borrow the bond.
\item \textbf{Cap: the GC rate.} If the special rate could not fall below 0\%, the special \emph{spread} cannot exceed the GC rate --- otherwise a lender would simply do a general collateral repo instead.
\item \textbf{Since 2009: the penalty rate.} Until 2009 there was no penalty for failed trades. After the crisis and trillions of dollars in failed OTR deliveries, regulators adopted a penalty rate equal to the \textbf{greater of 3\% minus the federal funds rate, or zero}. As the fed funds rate \emph{rises}, the penalty \emph{falls}; when the fed funds rate reaches zero, the penalty reaches its \textbf{maximum of 3\%}. The penalty rate becomes the new upper limit for the special spread.
\end{enumerate}
\end{lrkeybox}

\begin{lrtrapbox}[The penalty rate moves against the fed funds rate]
$\text{Penalty} = \max(3\% - \text{FFR},\; 0)$. That is an \textbf{inverse} relationship, and it is the single most corruptible fact in this objective. High rates $\rightarrow$ low penalty. Zero rates $\rightarrow$ the full 3\%.
\end{lrtrapbox}

% =====================================================================
\lo{LTR-13 g}{Calculate the financing advantage of a bond trading special when used in a repo transaction.}

\begin{lrkeybox}[What the source notes carry]
The financing value of the bond is the ability to \textbf{lend the bond at a relatively cheap special rate and use the cash to lend out at higher GC rates}. This financing value depends on the trader's expectation of \textbf{how long the bond will continue trading at its special rate}.
\end{lrkeybox}

\gapbadge
\gapnote{Neither layer contains the calculation. The Crash layer gives the description above and stops; the Lecture page carrying this objective's heading contains Covered Interest Parity instead. The computation below is written from the GARP curriculum (Tuckman and Serrat, Chapter 12, equation 14.3).}

\begin{lrgapbox}
The objective is \emph{calculate}. What is needed is the conversion of a special spread into a price value, then into a yield value --- and the fact that OTR richness has \emph{two} components, of which financing is usually the smaller.
\end{lrgapbox}

\begin{lrfmlbox}[Financing value of a bond trading special]
\[ \text{Value per 100 of cash} \;=\; 100 \times \frac{\text{Days} \times \text{Special spread}}{360} \]
\[ \text{Value per 100 face} \;=\; \text{Value per 100 of cash} \times \frac{\text{Price}}{100} \qquad
   \text{Yield value (bp)} \;=\; \frac{\text{Value per 100 face}}{\text{DV01}} \]
\vk{Days = the horizon over which the bond is expected to keep trading special; special spread = GC rate minus special rate, as a decimal; price = the bond's market price per 100 face; DV01 = the dollar value of a basis point per 100 face. The day count is \textbf{actual/360}, as elsewhere in repo.}
\end{lrfmlbox}

\begin{lrexbox}[Valuing the specialness of the OTR 10-year]
The on-the-run 10-year, the 3\textonehalf s of 15 May 2020, was trading extremely rich on 28 May 2010 --- slightly more than 2 per 100 face relative to the C-STRIPS curve. The term special spread to 30 September 2010 is \textbf{0.22\%} over \textbf{125 days}. The price is \textbf{101.90} and the DV01 is approximately \textbf{0.085}. Assume no haircuts. What is the financing advantage worth?
\tcblower
\textbf{Step 1 --- value of lending 100 of cash at the special spread for 125 days.}
\[ 100 \times \frac{125 \times 0.22\%}{360} = 100 \times \frac{0.275}{360} = \mathbf{0.076} \]
or \textbf{7.6 cents per 100 market value} of the bond.

\textbf{Step 2 --- convert from market value to face value} by scaling at the price of 101.90.
\[ 0.076 \times 1.019 = \mathbf{0.077} \]
or about \textbf{7.7 cents per 100 face amount}.

\textbf{Step 3 --- translate the dollar value into a yield value} by dividing by the DV01.
\[ \frac{0.077}{0.085} = \mathbf{0.9 \text{ basis points}} \]

\textbf{What this means, and why it is the point of the objective.} The bond was rich by \textbf{more than 2 dollars}, and its financing advantage accounts for \textbf{7.7 cents} of that --- under 4\%. OTR bonds trade at a premium partly for their \textbf{liquidity advantage} (the ability to turn positions back into cash with minimum effort, even in a crisis) and partly for their \textbf{financing advantage} (the ability to lend the bond and borrow cash cheaply). In this case almost all of the premium is due to \textbf{liquidity}, not financing. A candidate who assumes specialness explains OTR richness has the story backwards.
\end{lrexbox}

\begin{lrtrapbox}[Three failure modes in this calculation]
\begin{itemize}
\item \textbf{Stopping at step 1.} 7.6 cents is per 100 \emph{market value}; the answer per 100 \emph{face} is 7.7. Both will be offered, and the difference is one multiplication by the price.
\item \textbf{Day count.} Actual/360. Using 365 gives 0.0753 instead of 0.0764.
\item \textbf{Dividing the wrong way at step 3.} Dollar value $\div$ DV01 gives basis points. Multiplying instead gives 0.0065, which is not a plausible magnitude but reads as a number.
\end{itemize}
\end{lrtrapbox}

\begin{lrtrapbox}[Consolidated trap summary — LTR-13]
\begin{itemize}
\item \textbf{Input twin.} Repo settlement uses the \textbf{market value} contract price, not face value, on an \textbf{actual/360} basis. The general form carries \textbf{accrued interest} and a \textbf{haircut}; \S5B.9 records three distractors, one per omitted input.
\item \textbf{Polarity.} In repo, \textbf{counterparty risk is the smaller concern} (collateral mitigates it) and \textbf{liquidity risk is the bigger one}.
\item \textbf{Ordering.} FFR $>$ GC rate $>$ special rate. GC is the \emph{highest} repo rate; specials trade \emph{below} it.
\item \textbf{Polarity.} Special spreads are \textbf{narrower after an auction} (new supply) and \textbf{wider before} one (rising demand for the bond being borrowed).
\item \textbf{Polarity.} The failed-trade penalty is $\max(3\% - \text{FFR}, 0)$ --- it \textbf{falls} as rates rise and peaks at 3\% when rates are zero.
\item \textbf{Scope.} The special \emph{rate} floors at 0\%; the special \emph{spread} caps at the GC rate, and since 2009 at the penalty rate.
\item \textbf{Intermediate result.} 7.6 cents (per 100 market value) versus 7.7 cents (per 100 face) versus 0.9 bp --- three correct answers to three different questions in one chain.
\item \textbf{Scope.} OTR richness is mostly \textbf{liquidity}, not financing. In the worked case, 7.7 cents of a premium above 2 dollars.
\item \textbf{Role.} Lehman: intraday haircuts from its own tri-party clearing agent. Bear Stearns: a three-sided run in which even agency collateral could not buy a rollover.
\end{itemize}
\end{lrtrapbox}
